\section{The functional limit and its randomness}
\label{sec:functional}

Proposition~\ref{prop:trace-goe} and
Theorem~\ref{thm:comparison-universality} identify the same finite
collections of resolvent traces for Gaussian and sign disorder.
We first turn these statements into convergence of autocorrelation
functions.  We then prove decay and show that the edge dependence
produces a random, rather than deterministic, limit.

\subsection{From resolvents to continuous functions}

Let $E_n$ be the spectral resolution of the nonnegative operator $TH$,
and define the finite measure
\[
 \nu_n(I)=\sum_{i=1}^n
 \langle n^{-1/3}X_i,E_n(I)n^{-1/3}X_i\rangle_\mu
\]
for Borel sets $I\subset[0,\infty)$.  Basis invariance gives
\[
 B_n(z)=\int\frac{d\nu_n(\lambda)}{z+\lambda},
 \qquad
 A_{n,J}(s)=\int e^{-s\lambda}\,d\nu_n(\lambda).
\]
Although $\nu_n([0,\infty))=n^{1/3}$ diverges, its division by
$1+\lambda$ has tight mass.  This is the reason to use resolvents
rather than the unweighted spectral measures.

\begin{lemma}[A continuity principle for spectral measures]
\label{lem:functional-spectral}
Let $\nu_n$ and $\nu$ be random nonnegative measures on $[0,\infty)$
such that their integrals against $(1+\lambda)^{-1}$ are finite almost
surely, for every $n$ and for $\nu$.
Suppose the corresponding integrals against $(z+\lambda)^{-1}$
converge jointly in law for every finite collection of positive $z$,
including $z=1$.  Then
\[
 \left(s\longmapsto\int e^{-s\lambda}\,d\nu_n(\lambda)\right)
 \ \Longrightarrow\
 \left(s\longmapsto\int e^{-s\lambda}\,d\nu(\lambda)\right)
\]
in $C_{\rm loc}((0,\infty))$.
\end{lemma}

\begin{proof}
Set $d\eta_n=(1+\lambda)^{-1}d\nu_n$ and regard $\eta_n$ as a finite
measure on the compact space $K=[0,\infty]$, giving no mass to infinity.
Its total mass is the Stieltjes transform at $z=1$, and is therefore
tight.  Finite measures on $K$ with bounded total mass form a compact
set, so the laws of $\eta_n$ are tight.

For $z>0$ the function
\[
 k_z(\lambda)=\frac{1+\lambda}{z+\lambda},\qquad k_z(\infty)=1,
\]
is continuous on $K$.  A countable dense family of these functions,
including $k_1=1$, determines a finite measure even when that measure
has mass at infinity.  Indeed their closed linear span contains
$1$ and $r_z(\lambda)=(z+\lambda)^{-1}$, with $r_z(\infty)=0$.
The partial-fraction identity
$r_zr_w=(r_z-r_w)/(w-z)$, followed by uniform limits when $w=z$,
shows that this span is an algebra.  It separates the points of $K$,
so it is dense in $C(K)$.  Consequently every subsequential law is
the law of $d\eta=(1+\lambda)^{-1}d\nu$, with no mass at infinity.
Here joint convergence on the countable determining family identifies
the law of the random measure, rather than only its mean.

For $s>0$, define $h_s(\lambda)=(1+\lambda)e^{-s\lambda}$ and
$h_s(\infty)=0$.  On each compact interval $a\leq s\leq b$, $a>0$,
the map $s\mapsto h_s$ is continuous in the uniform norm on $K$.
Thus weak convergence of finite measures on $K$ gives convergence
of their $h_s$ integrals uniformly on $[a,b]$: a finite uniform-norm
net of the functions $h_s$ reduces the assertion to finitely many
continuous tests, while weak convergence bounds the total masses.
The resulting map from measures to $C_{\rm loc}((0,\infty))$ is
continuous.  The continuous mapping theorem proves the claim.
\end{proof}

\begin{proposition}[Identification of the functional limit]
\label{prop:functional-identification}
The series \eqref{eq:source-2} has a Borel version
$g\mapsto\mathcal A_g$ with values in $C_{\rm loc}((0,\infty))$.
For both Gaussian and symmetric-sign entries,
\[
 \operatorname{Law}(A_{n,J})\Longrightarrow
 Q_{\rm crit}=\operatorname{Law}(\mathcal A_g),
\]
where the same edge law and the same deterministic coefficient $c_*$
are used in \eqref{eq:source-2}.
\end{proposition}

\begin{proof}
For a fixed edge configuration $g$, let $E_g$ be the spectral resolution
of $c_*H_\Pi$ and put
\[
 \nu_g(I)=\sum_{a\geq1}\langle x_a,E_g(I)x_a\rangle_\Pi.
\]
This countable sum defines a nonnegative measure.  By
Proposition~\ref{prop:trace-goe}, its division by $1+\lambda$ has
finite mass $B_g(1)$ almost surely.  It follows in particular that
the series in \eqref{eq:source-2} defines a continuous function
$\mathcal A_g$ on $(0,\infty)$.

These constructions are measurable in $g$.  The density prescription
for $\Pi$ makes each cylinder integral measurable.  Resolvent pairings
can then be computed by the variational supremum over a countable
family of smooth compact cylinders.  The family can be chosen
independently of $g$ by ordinary approximation of a function and its
first derivatives on finite-dimensional compact sets; it is a form
core by Proposition~\ref{prop:gradient-domain}.  Thus the resolvent
pairings, the divided spectral measure determined by them, and
$\mathcal A_g$ are measurable.

Apply Lemma~\ref{lem:functional-spectral} to the finite-dimensional
trace convergence in Proposition~\ref{prop:trace-goe} and
Theorem~\ref{thm:comparison-universality}.  For either entry law,
\[
 \operatorname{Law}(A_{n,J})\Longrightarrow
 Q_{\rm crit}:=\operatorname{Law}(\mathcal A_g)
 \quad\text{in }C_{\rm loc}((0,\infty)).
\]
The common coefficient is the deterministic $c_*$ of
Theorem~\ref{thm:micro-relaxation}; no further change of time is involved.
\end{proof}

\subsection{Decay and positive values}

The spectral representation makes $\mathcal A_g$ nonnegative and
nonincreasing.  Proposition~\ref{prop:gradient-domain} identifies the
kernel of $H_\Pi$ as the constants.  Each $x_a$ is centered by sign
symmetry of $\Pi$, so $\nu_g$ has no atom at zero.  For any fixed
$s_0>0$, $e^{-s_0\lambda}$ is integrable against $\nu_g$.
Dominated convergence for $s\geq s_0$ therefore gives
\[
 \lim_{s\to\infty}\mathcal A_g(s)=0.
\]

There is also a useful positive lower bound.  Put
$v(g)=\Pi[x_1^2]$, which is finite by the coordinate moment bounds and
strictly positive by full support of the first-coordinate marginal.
The coordinate is
in the form domain with energy one, by
Proposition~\ref{prop:gradient-domain}.  Jensen's inequality for its
normalized spectral measure gives
\begin{equation}
 \mathcal A_g(s)\geq
 \langle x_1,e^{-s c_*H_\Pi}x_1\rangle_\Pi
 \geq v(g)\exp\{-s c_*/v(g)\}>0.
 \label{eq:functional-positive}
\end{equation}
It remains to show that these functions do not have a deterministic law.
For this we will vary a leading edge point while retaining absolute
continuity with respect to the original edge law.

\subsection{Finite shifts of the edge process}

Let $\mathsf P$ denote the law of the increasing edge sequence $g$.
For an integer $m$ and a vector $h\in\mathbb R^m$, define
\[
 T_hg=(g_1+h_1,\ldots,g_m+h_m,g_{m+1},\ldots),
\]
and let $\mathcal O_h$ be the event that this sequence is still strictly
increasing.  We use the restriction of $\mathsf P$ to this event,
without normalizing its mass.

\begin{lemma}[Absolute continuity under finite shifts]
\label{lem:functional-shift}
For every fixed $m$ and $h\in\mathbb R^m$,
\[
 (T_h)_*(\mathsf P|_{\mathcal O_h})\ll\mathsf P.
\]
In particular, lowering $g_1$ by any fixed positive amount gives a
law absolutely continuous with respect to $\mathsf P$.
\end{lemma}

\begin{proof}
We prove the claim by finite-dimensional GOE density ratios.  Let
$\mathsf P_n$ denote the ordered GOE edge law in dimension $n$.
When a sequence-valued argument is needed, append the values
$g_{n,n}+j$, $j\geq1$, after its final coordinate.  This convention
does not affect any fixed coordinate for large $n$, so
$\mathsf P_n\Longrightarrow\mathsf P$ in the product topology; it
also commutes with $T_h$ whenever $n>m$.
In the
variables $g_{a,n}=n^{2/3}(2-\lambda_a)$, the ordered density is
proportional to
\[
 \exp\left\{-\frac n4\sum_{a=1}^n
                 (2-n^{-2/3}g_{a,n})^2\right\}
 \prod_{i<j}(g_{j,n}-g_{i,n}).
\]
Translation of the first $m$ variables has Jacobian one.  Its Gaussian
log-density ratio is
$n^{1/3}\sum_{k\leq m}h_k+O(1)$ on sets where the first $m$
coordinates are bounded.  For a fixed $k\leq m$, the spectral estimates
of Section~\ref{sec:edge-law} give, tightly,
\[
 \sum_{m<i\leq n}\frac1{g_{i,n}-g_{k,n}}=n^{1/3}+O(1),
 \qquad
 \sum_{m<i\leq n}\frac1{(g_{i,n}-g_{k,n})^2}=O(1).
\]
To see the first assertion, compare the reciprocals with
$(g_{i,n}+b_n)^{-1}$, where $g_{1,n}+b_n\geq1$.
Their difference is summable by the square-tail bound, while
$n^{1/3}-\sum_i(g_{i,n}+b_n)^{-1}$ is tight.  Tight separation of
the first finitely many eigenvalues controls the remaining denominators.

The logarithm of the Vandermonde ratio involving indices $k\leq m<i$
contains
\[
 \sum_{m<i\leq n}\log\left(1-\frac{h_k}{g_{i,n}-g_{k,n}}\right).
\]
Beyond a sufficiently large fixed index, the fractions have absolute
value at most $1/2$.  Taylor's inequality and the preceding two bounds
show that this sum is $-h_kn^{1/3}+O(1)$.  The finitely many remaining
factors, including pairs within the first $m$ coordinates, are bounded
away from zero after imposing positive margins on the old and shifted
gaps.  The linear terms cancel those of the Gaussian factor.
Consequently the new-to-old density ratio has a positive lower bound
on sets of probability arbitrarily close to one, after any fixed
localization inside the strict-order event.

Here is the measure-theoretic direction of the resulting implication.
Let $\chi$ be a continuous function of finitely many leading coordinates,
with $0\leq\chi\leq1$, supported where the shifted order has a positive
margin and the leading coordinates are bounded.  For every
$\epsilon>0$, the density-ratio bound supplies $c_\epsilon>0$ such that,
for every bounded nonnegative continuous function $f$ of the edge
sequence,
\[
 \mathsf P_n(f)\geq c_\epsilon
 \left(\mathsf P_n[\chi\,f\circ T_h]
                       -\epsilon\|f\|_\infty\right).
\]
The common spectral exceptional sets have probability at most
$\epsilon$ for all sufficiently large $n$.  Pass to the edge limit.
For a compact subset of a $\mathsf P$-null set, take continuous tests
between zero and one, equal to one on that compact subset and supported
in open sets of arbitrarily small $\mathsf P$-mass.  The inequality
bounds its $(T_h)_*(\chi\mathsf P)$-mass by $\epsilon$.
Inner regularity gives the same bound for every $\mathsf P$-null set.
Letting $\epsilon\downarrow0$ proves
absolute continuity for this localized measure.  Finally exhaust
$\mathcal O_h$ by such continuous localizations.  This proves the
lemma without requiring a uniform bound on the density ratio over
the entire edge law.
\end{proof}

\begin{proposition}[Randomness of the limit]
\label{prop:functional-randomness}
For every fixed $s>0$, the finite random variable $\mathcal A_g(s)$
is essentially unbounded.  In particular $Q_{\rm crit}$ is not a
point mass.
\end{proposition}

\begin{proof}
Lower $g_1$ by $t>0$, leaving the other points fixed, and denote the
resulting sequence by $g^{(t)}$.  Choose the auxiliary parameter $b$
large enough for both Gaussian prescriptions.  Only $l_1$ changes,
and \eqref{eq:source-1} shows
\[
 D_{g^{(t)}}(b)-D_g(b)
 =-\big(l_1^{(t)}-l_1\big).
\]
At a common coordinate sequence $x$, the finite centered-square sums
in the two prescriptions differ by $-(l_1^{(t)}-l_1)$.  Thus, wherever
these sums converge,
$S_{g^{(t)}}(x)=S_g(x)-(l_1^{(t)}-l_1)$.  Consequently
\[
 S_{g^{(t)}}(x)-D_{g^{(t)}}(b)=S_g(x)-D_g(b).
\]
The two renormalized constraints therefore coincide.  Comparing the
first Gaussian coordinate in their
conditioning densities gives
\[
 \frac{d\Pi_{g^{(t)}}}{d\Pi_g}(x)
 =\frac{\exp(t x_1^2/2)}{\Pi_g[\exp(t x_1^2/2)]}.
\]
The denominator is finite: increase $b$ until $g_1+b>t$ and use the
bounded conditional first-coordinate density relative to the centered
Gaussian of variance $(g_1+b)^{-1}$.  The independence of the
conditioned law from $b$ then gives the displayed identity for $\Pi_g$.

For a nonnegative random variable with unbounded support, exponential
tilts with parameter tending to infinity have means tending to
infinity.  Indeed, for any $M$, the mass below $M$ after tilting is
at most
$e^{-t/2}\Pi_g[x_1^2\leq M]/\Pi_g[x_1^2\geq M+1]$.
It follows that $v(g^{(t)})\to\infty$ as $t\to\infty$.

Suppose that $v(g)$ had an essential finite upper bound.  By
Lemma~\ref{lem:functional-shift}, this same bound would hold for
$v(g^{(t)})$ almost surely for every positive integer $t$.
Intersecting these countably many full-measure events contradicts
the preceding divergence.  Therefore $v(g)$ is essentially unbounded.
The lower bound \eqref{eq:functional-positive} proves the same assertion
for $\mathcal A_g(s)$ at every fixed positive $s$.
\end{proof}

\begin{proof}[Completion of the proof of Theorem~\ref{thm:main}]
Proposition~\ref{prop:functional-identification} proves convergence of the laws in
$C_{\rm loc}((0,\infty))$ for both entry distributions and identifies
their common limit by \eqref{eq:source-2}.  The same deterministic
$c_*>0$ occurs in both models.  The spectral representation and the
trivial form kernel prove nonnegativity, monotonicity, and almost-sure
decay.  Equation~\eqref{eq:functional-positive} excludes the zero
function, and Proposition~\ref{prop:functional-randomness} establishes
sample-to-sample randomness.
\end{proof}

Decay at infinity is not a closed condition in
$C_{\rm loc}((0,\infty))$.  The almost-sure decay proved above can
therefore coexist with nondecaying functions in the closed topological
support.  Appendix~\ref{sec:support} gives a precise example: the
constant-one function belongs to the closed topological support of
$Q_{\rm crit}$.
